\section{Axiom 1: Element-Wise Independence}
\label{sec:axiom-1}

% Cross-references to sec:miller-rabin (ch 3) and ch:approximate-sets (ch 5)
% will resolve when those chapters are drafted in later plans.
% Build warnings for those labels are expected and will be resolved in Plan 4.

The previous section sketched two kinds of independence that chapters 1--3
quietly assumed. This section makes the first of them precise.

A brief note on framing before we begin. In \Cref{sec:why-axioms} we
used a chain-of-stages example as motivation: $k$ noisy channels in
series, with end-to-end error rate $1 - \prod_{i=1}^k (1 - \tnrate_i)$.
That example invokes a kind of independence, but it is not the same
statement as Axiom 1. Axiom 1 governs independence across distinct input
elements of a \emph{single} approximate predicate. Stage-wise independence
in a composition chain is a separate assumption: closely related, often
co-occurring, but operationally distinct. Both are needed for the chain
formula; both fail in similar circumstances; and in well-behaved
implementations both hold. We state Axiom 1 cleanly here, on a single
predicate. The reader who wants to recover the chain formula from these
axioms should track the two independences separately and combine them as
the composition theorem of \Cref{ch:composition} prescribes.

\subsection*{The formal statement}

Let $U$ be the universal set, the domain of all possible inputs to
the approximate membership predicate. For a fixed objective set
$A \subseteq U$, the predicate either correctly reports the true answer
or makes an error. For each $x \in U$, define the error indicator
\[
  E_x \;=\; \mathbf{1}\bigl[\tilde{A}(x) \neq \mathbf{1}_A(x)\bigr],
\]
where $\tilde{A}(x)$ is the approximate predicate's output on input $x$
and $\mathbf{1}_A(x)$ is the exact answer. $E_x = 1$ means an error
occurred at $x$; $E_x = 0$ means the answer was correct.

\begin{axiom}[Element-wise independence]
\label{ax:element-indep}
For all distinct $x, y \in U$, the error events $E_x$ and $E_y$ are
independent:
\[
  \Prob{E_x = 1, \, E_y = 1} \;=\; \Prob{E_x = 1} \cdot \Prob{E_y = 1}.
\]
More generally, the family $\{E_x : x \in U\}$ is mutually independent
given the objective set $A$ and the error rates. Furthermore, within
each partition block $B_i$ of $U$ (defined formally in \Cref{sec:axiom-2}),
the error indicators $\{E_x : x \in B_i\}$ are identically distributed.
\end{axiom}

The notation $\tilde{A}(x)$ is deliberately abstract: it stands for any
approximate membership test, whether implemented as a Bloom filter, a
probabilistic data structure, a noisy sensor, or a randomized algorithm.
Axiom 1 is a property of the \emph{process} that generates the answers,
not of any particular implementation.

\subsection*{Consequence: factored joint distributions}

Axiom 1 is not merely a philosophical assertion. It has an immediate
and quantitative consequence for joint query distributions.

\begin{proposition}[Factored joint distribution]
\label{prop:factored-joint}
Under \Cref{ax:element-indep}, the joint distribution of the error
indicators over any finite set of distinct inputs $\{x_1, \ldots, x_n\}
\subseteq U$ factors as a product of $n$ marginal distributions:
\[
  \Prob{E_{x_1} = e_1, \, \ldots, \, E_{x_n} = e_n}
  \;=\;
  \prod_{i=1}^n \Prob{E_{x_i} = e_i}
\]
for any choice of $e_1, \ldots, e_n \in \{0, 1\}$.
Equivalently, the joint distribution of the observed responses
$\tilde{A}(x_1), \ldots, \tilde{A}(x_n)$ factors as
\[
  \Prob{\tilde{A}(x_1) = s_1, \ldots, \tilde{A}(x_n) = s_n \Given A}
  \;=\;
  \prod_{i=1}^n \Prob{\tilde{A}(x_i) = s_i \Given A},
\]
where each factor equals $\fprate$, $\tnrate$, $\tprate$, or $\fnrate$
depending on whether $x_i \in A$ and whether $s_i$ is a correct response.
\end{proposition}

\begin{proof}
The first identity is the definition of mutual independence for the
family $\{E_x : x \in U\}$, which Axiom 1 asserts directly.
The second identity follows because $\tilde{A}(x_i) = s_i$ is a
deterministic function of $E_{x_i}$ and the exact answer $\mathbf{1}_A(x_i)$:
knowing $A$ determines $\mathbf{1}_A(x_i)$, and given that,
$\tilde{A}(x_i) = s_i$ iff $E_{x_i}$ takes the appropriate value.
The product therefore inherits its factored form from the independence
of $\{E_{x_i}\}$.
\end{proof}

This proposition is the engine behind every multi-query formula in the
book. Whenever we write a joint probability over several distinct inputs
as a product of marginals, we are using this proposition, which rests
on Axiom 1.

The factored joint distribution also licenses a key observation about
the joint confusion matrix of $m$ simultaneous queries. Without any
independence assumption, that matrix lives in a space of dimension
$n(2^m - 1)$ for an $n$-th order model. Axiom 1 forces it to be a
Kronecker product of per-query $2 \times 2$ matrices, collapsing the
parameter count to at most $2m$. That is the content of
\Cref{sec:kronecker}; we state it here as a forward reference so the
reader knows why Axiom 1 matters at the level of the whole framework.

\subsection*{Example 1: The Bloom filter's independence assumption}

A Bloom filter~\cite{bloomfilter} answers membership queries on a set
$A$ stored in a bit array of $m$ bits, using $k$ independent hash
functions. When $k$ hashes are applied to a new element $x$, each hash
selects a bit position; $x$ is declared a member if and only if all $k$
selected positions are 1.

For elements $x \notin A$, the probability that all $k$ positions happen
to be set (a false positive) is approximately $(1 - e^{-kn/m})^k$, where
$n = |A|$. This familiar formula is derived under an explicit assumption:
for two distinct non-members $x$ and $y$, the event ``$x$ is a false
positive'' is independent of the event ``$y$ is a false positive.'' Without
this assumption, the formula gives neither an exact probability nor a
useful bound; instead, one would need to account for correlations between
all pairs of non-members, producing a quantity with no closed form.

Axiom 1 is exactly the assumption the Bloom filter's analysis needs.
The filter's $k$ hash functions are chosen to be independent, and the
hash values for distinct elements are treated as unrelated. Whether or
not this is precisely true in a given implementation, it is the
assumption under which the false-positive rate formula is derived, and
it is the assumption that the Bernoulli framework formalizes.

We will analyze the Bloom filter fully in Part II
(\Cref{ch:approximate-sets}), where Axiom 1 will be verified to hold
under the standard random-oracle model for the hash functions. For now,
the take-away is structural: the Bloom filter is a $\Bernoulli{\BoolType}$
implementation, and Axiom 1 is the assumption that makes that claim
non-trivial.

\subsection*{Example 2: Miller-Rabin's $k$-round bound requires Axiom 1}

In \Cref{sec:miller-rabin} we derived the Miller--Rabin bound: after
$k$ independent rounds of witness testing, the probability that a
composite number passes all $k$ rounds (the $k$-round false-positive
rate) is at most $(1/4)^k$. The bound rests on two facts: each round's
false-positive probability is at most $1/4$, and the rounds are independent.

There is a subtle reframing required to fit Miller-Rabin under Axiom 1.
The predicate being tested takes a \emph{single} fixed input $n$ (the
integer in question), and randomness enters via the witness draw rather
than via the input. To apply Axiom 1 cleanly, recast the test as a
predicate $\tilde{P}\colon \{2, \ldots, n-2\} \to \{\text{composite},
\text{probably prime}\}$ on the witness pool: now distinct witnesses
\emph{are} distinct input elements, and Axiom 1's claim is that error
events at distinct witnesses (the events ``$w_i$ is a strong liar for
$n$'') are independent. Whether or not this independence holds for any
particular composite $n$ depends on the structure of $n$'s strong-liar
set, a separate fact about Miller-Rabin's mathematical content.

With that reframing in place, the connection to Axiom 1 is clean. Each
round tests a different witness $w_i$ drawn uniformly at random from
$\{2, \ldots, n-2\}$. Whether $w_i$ happens to be a strong liar for $n$
(the error event at $w_i$) must be independent of whether $w_j$ is a
strong liar (the error event at $w_j$), for $i \neq j$. The error events
are the events that the witness fails to detect compositeness.

Under Axiom 1, the joint probability that all $k$ witnesses are strong
liars is the product of the individual probabilities:
\[
  \Prob{\text{all } k \text{ witnesses are strong liars}}
  \;\leq\;
  \prod_{i=1}^k \frac{1}{4} \;=\; \left(\frac{1}{4}\right)^k.
\]
Without Axiom 1, the $k$ error events could be positively correlated.
The events ``$w_1$ is a liar'' and ``$w_2$ is a liar'' may be
positively correlated, because both witnesses are sampled from the same
composite $n$, which may itself happen to have an unusually large
strong-liar set. Conditioning on $n$, the witnesses are not independent
of each other when the liar fraction is unknown. In that
case, the product bound $(1/4)^k$ would underestimate the true
false-positive probability. The correct expression would involve the
joint distribution over all $k$ witnesses, which does not simplify to
a product without an independence assumption.

The Miller--Rabin analysis works because the witnesses are drawn
uniformly and independently from a large set. In the worst case, at
most $1/4$ of all candidates are strong liars; with independent uniform
sampling, the fraction is driven down exponentially.
This is Axiom 1 at work.

\subsection*{Classifier-verdict callback: a related independence failure}

The classifier-verdict thread from \Cref{sec:why-axioms} returns here,
but with a deliberate caveat. Two correlated classifiers evaluating the
\emph{same} input is not literally an Axiom 1 failure: Axiom 1 governs
distinct \emph{inputs} to a single predicate, not distinct
\emph{predicates} on a single input. The example we work below is one
step removed from Axiom 1, but it illustrates the same family of
phenomenon: when error events are correlated for any reason, the naive
product formula underestimates the joint failure rate. The Axiom 1 case
(correlated errors at distinct inputs to one predicate) is
mathematically structurally identical.

Suppose two classifiers $C_1$ and $C_2$ are both trained on a shared
dataset with overlapping features. Each classifier is a
$\Bernoulli{\BoolType}$ individually: on any given input $x$,
$C_1$ has false-positive rate $\fprate_1$ and $C_2$ has false-positive
rate $\fprate_2$.

Now consider their joint behavior on a single negative input $x \notin A$.
If the errors were independent, the probability that both make a false
positive would be:
\[
  \Prob{C_1(x) = 1, \, C_2(x) = 1 \Given x \notin A}
  \;\stackrel{?}{=}\;
  \fprate_1 \cdot \fprate_2.
\]
But the two classifiers share training data, so their errors on $x$
tend to be correlated.

Let $\rho \in [0,1]$ be the correlation between the error indicators
$E_x^{(1)}$ and $E_x^{(2)}$. The joint probability is
\[
  \Prob{E_x^{(1)} = 1, \, E_x^{(2)} = 1}
  \;=\;
  \fprate_1 \fprate_2 \;+\; \rho \sqrt{\fprate_1(1-\fprate_1) \cdot \fprate_2(1-\fprate_2)}.
\]
For a concrete example, let $\fprate_1 = \fprate_2 = 0.10$ and
$\rho = 0.60$. The product bound predicts a joint false-positive rate of
$0.10 \times 0.10 = 0.01$. The actual joint rate is
\[
  0.01 + 0.60 \times \sqrt{0.10 \times 0.90 \times 0.10 \times 0.90}
  \;=\;
  0.01 + 0.60 \times 0.09
  \;=\;
  0.01 + 0.054
  \;=\;
  0.064.
\]
The inflation factor is $0.064 / 0.01 = 6.4$: the true joint
false-positive rate is $6.4\times$ the naive product.

This is not a failure of the individual classifiers, which are each
well-behaved $\Bernoulli{\BoolType}$ objects; it is a failure of the
independence assumption between them. The Axiom 1 violation has the same
arithmetic; only the source of the correlation differs.

The twin lessons are these. First, Axiom 1 is a genuine modeling
assumption, not an automatic consequence of each individual predicate
being a $\Bernoulli{\BoolType}$: the individual marginals can be well
calibrated while the joint distribution fails to factor. Second,
correlated errors are worse than independent errors in a precise,
quantifiable sense: the inflation factor grows with the correlation
coefficient and with the individual error rates.

The Kronecker theorem proved in \Cref{sec:kronecker} makes this
precise at the level of the joint confusion matrix: under Axiom 1,
the matrix has a structured Kronecker-product form that encodes exactly
the factored joint distribution in \Cref{prop:factored-joint}. Without
Axiom 1, the matrix is unstructured, the parameter count explodes, and
the elegant formulas of chapters 1--3 must be replaced by unwieldy
joint distributions over all pairs and triples of query events.
